\section{QCD gradient flow in perturbation theory}\label{sec:pert}

In the following, we will work in $D$-dimensional Euclidean space-time
with $D=4-2\ep$.  The gradient-flow formalism continues the gluon and
quark fields $A^a_\mu(x)$ and $\psi(x)$ of regular\footnote{We will use
  the terms ``flowed'' and ``regular'' \qcd\ to distinguish quantities
  defined at $t>0$ from those defined at $t=0$.} \qcd\ to
$(D+1)$-dimensional fields $B^a_\mu(t,x)$ and $\chi(t,x$) through the
boundary conditions
\begin{equation}
  \begin{split}
    B_\mu^a (t=0,x) = A_\mu^a (x)\,,\qquad \chi (t=0,x)= \psi(x)\,,
    \label{eq:bound}
  \end{split}
\end{equation}
and the flow equations
\begin{equation}
  \begin{split}
    \partial_t B_\mu^a &= \mathcal{D}^{ab}_\nu G_{\nu\mu}^b + \kappa
    \mathcal{D}^{ab}_\mu \partial_\nu B_\nu^b\,,\\ \partial_t \chi &= \Delta \chi
    - \kappa \partial_\mu B_\mu^a T^a \chi\,,\\ \partial_t \overline
    \chi &= \overline \chi \overleftarrow \Delta + \kappa \overline \chi
    \partial_\mu B_\mu^a T^a\,,
    \label{eq:flow}
  \end{split}
\end{equation}
where the ``flow time'' $t$ is a parameter of mass dimension minus two, and
$\kappa$ is an additional gauge parameter which drops out of physical
observables.

The $(D+1)$-dimensional field-strength tensor is defined as
\begin{align}
  G_{\mu\nu}^a = \partial_\mu B_\nu^a - \partial_\nu B_\mu^a + f^{abc}
  B_\mu^b B_\nu^c\,,
\end{align}
the covariant derivative in the adjoint representation is given by
\begin{align}
  \mathcal{D}_\mu^{ab} = \delta^{ab} \partial_\mu - f^{abc} B_\mu^c\,,
\end{align}
and
\begin{equation}
  \begin{split}
    \Delta = (\partial_\mu + B_\mu) (\partial_\mu + B_\mu)\,,\qquad
    \overleftarrow{\Delta} = (\overleftarrow\partial\!_\mu - B_\mu)
    (\overleftarrow\partial\!_\mu - B_\mu)\,.
  \end{split}
\end{equation}
As usual, the color indices of the adjoint representation are denoted by
$a,b,c,\ldots$, while $\mu,\nu,\rho,\ldots$ are $D$-dimensional Lorentz
indices. Color indices of the fundamental representation are suppressed
throughout this paper, unless required by clarity.  The symmetry
generators $T^a$ obey the commutation relation
\begin{equation}
  \begin{split}
    [T^a,T^b] = f^{abc}T^c\,,
  \end{split}
\end{equation}
with the structure constants $f^{abc}$.

The flow-field equation leads to a smearing of gauge-field
configurations at finite flow time $t>0$. As a consequence, composite
operators at $t>0$ do not require renormalization beyond the
renormalization of the parameters and fields of the Lagrangian. For the
strong coupling and the quark mass, the renormalization constants are
identical to those at $t=0$; the flowed gluon fields do not require
renormalization at finite flow time as was pointed out in \citere{Luscher:2011bx}.
The renormalization constant for the flowed quark field through
\nnlo\ is a by-product of this paper and will be given below.
